\section{准分布（Quasi-Distributions）}

\subsection{定义及与 TMDs 的关系}

由于格点上不可能有类光间隔，
文献 \cite{Ji:2013dva} 建议考虑
类空间隔
$z= (0,0,0,z_3)$ [为简便记 \mbox{$z=z_3$}]。于是，在
\mbox{$p=(E, 0_\perp, P)$} 参考系中，
通过如下参数化引入 quasi-PDF $Q(y, P)$：
\begin{align}
\langle p |   \phi(0) \phi (z_3)|p \rangle
=  &
\int_{-\infty}^{\infty}   dy \,
Q(y, P) \,  e^{i y  P z_3 } \,
\  .
\label{newVDFxzQ}
\end{align}
按照这一定义，函数 $Q(y,p)$ 刻画部分子携带强子第三动量分量 $P$ 的份额为 $y$
的概率。把该矩阵元素视为 $\nu$ 与 $-z^2$ 两个变量的函数
（在此参考系中它们由 $Pz_3$ 和 $z_3^2$ 给出），我们有
\begin{align}
{\cal M} (\nu,z_3^2)
=  &
\int_{-\infty}^{\infty}   dy \,
Q(y, P) \,  e^{i y  \nu } \,
\  .
\end{align}
注意到 $z_3^2= \nu^2/P^2$，便得到逆 Fourier 变换
\begin{align}
Q(y,  P)   =\frac{1}{2 \pi}  \int_{-\infty}^{\infty}  d\nu \,
\, e^{-i y  \nu}  \, {\cal M} (\nu,  \nu^2/P^2)    \  .
\label{IxM}
\end{align}
这表明，只要 ${\cal M} (\nu,  \nu^2/P^2) \to {\cal M} (\nu, 0)$，
在 $P \to \infty$ 极限下 $Q(y,  P)$ 趋于 $f(y)$。

因此，quasi-PDF $Q(y,  P)$ 偏离 PDF $f(y)$ 的程度
由 ${\cal M} (\nu,  z_3^2)$ 对其第二个宗量的依赖性决定。
根据 \mbox{式 (\ref{MTMD})，} 这一依赖性又与 TMD
${\cal F} (x, \kappa^2)$ 对其第二个宗量 $\kappa^2$ 的依赖性相关。
于是，$Q(y,  P)$ 与 $f(y)$ 之间的差别
可以用 TMDs 来描述。

为此，我们利用这样一个事实：\mbox{式 (\ref{MTMD})}
是函数 \mbox{${\cal M} (\nu,  z_1^2+z_2^2) $} 与
函数 ${\cal F} (x, k_1^2+k_2^2)$ 之间的一个{\it 数学}关系，
而不管变量
$z_1,z_2$ 与 $k_1,k_2$ 有什么{\it 物理}含义。
这样，把取 $z_1=0$、$z_2 =\nu/P$ 的 \mbox{式 (\ref{MTMD})}
代入式 (\ref{IxM})，
就把它转化为用 TMDs 表示 quasi-PDFs 的表达式：
\begin{align}
Q(y, P) /P =  & \,\int_{-\infty}^{\infty} d  k_1
\int_{-1} ^ {1} d  x \,   {\cal F} (x, k_1^2+(y-x)^2P^2 )
\  .
\label{QTMD}
\end{align}

最初，这一关系是在文献 \cite{Radyushkin:2016hsy} 中利用
文献 \cite{Radyushkin:2014vla,Radyushkin:2015gpa} 的 Nakanishi 型表示导出的。
现在我们看到，它仅仅是 Lorentz 不变性的一个直接推论。

\subsection{量子色动力学（QCD）情形}

上面导出的公式
可直接应用于 QCD 中非单态部分子密度。
此时处理的是形如
\begin{align}
{\cal M}^\alpha  (z,p) \equiv \langle  p |  \bar \psi (0) \,
\gamma^\alpha \,  { \hat E} (0,z; A) \psi (z) | p \rangle \
\label{Malpha}
\end{align}
的矩阵元素，其中 $
{ \hat E}(0,z; A)$ 是夸克（基础）表示中标准的 $0\to z$
直线规范链接。
这些矩阵元素可以分解为 $p^\alpha$ 与 $z^\alpha$ 两部分：
\begin{align}
{\cal M}^\alpha  (z,p) = &2 p^\alpha  {\cal M}_p (-(zp), -z^2)
\nn & + z^\alpha  {\cal M}_z (-(zp),-z^2)
\ .
\end{align}
当 $z^2 \to 0$ 时，${\cal M}_p (-(zp), -z^2) $ 部分给出扭度-2 分布，
而 $ {\cal M}_z ((zp),-z^2)$ 是纯高扭度污染，
最好把它去掉。

若在 ${\cal M}^\alpha$ 的 $\alpha=+$ 分量中取 $z=(z_-, z_\perp)$，
则 $z^\alpha$ 部分自动消失，
从而可以引入一个 TMD ${\cal F}(x, k_\perp^2)$，
它与 $ {\cal M}_p (\nu, z_\perp^2)$ 通过标量公式
(\ref{MTMD}) 相联系。
对于准分布，
去掉 $z^\alpha$ 污染的最简单办法是取
\mbox{$ {\cal M}^\alpha  (z=z_3,p)$} 的时间分量并定义
\begin{align}
&  {\cal M}^0   (z_3,p)
=  2p^0
\int_{-1}^1 dy\,
Q(y,P) \,
\,  e^{i  y Pz_3 }  \  .
\label{OPhixspin12}
\end{align}
于是 $Q(y,P)$ 与
${\cal F}(x, k_\perp^2)$ 之间的联系
由标量公式 (\ref{QTMD}) 给出。

可以注意到，定义 $ {\cal M}^\alpha  (z,p) $ 的算符
包含从 $0$ 到 $z$ 的直线链接，而非描述
Drell-Yan 与半单举 DIS 过程时所用 TMD 定义中常见的订书钉（staple）型链接。
众所周知，订书钉型链接
反映了这些过程中固有的初末态相互作用。
在这个意义上，``直线链接'' TMDs 描述的是强子处于
未受扰动或 ``原始'' 状态时的结构。虽然这样的 TMD
不太可能在散射实验中被测量到，
但它是一个良定义的 QFT 对象，
并且人们有希望能在格点上测量它。

\subsection{动量分布}

quasi-PDFs 描述的是部分子动量第三分量 $k_3$ 占强子动量份额
$y \equiv k_3/P$ 的分布。
也可以直接对 $k_3$ 本身引入分布：
$R (k_3,P) \equiv Q(k_3/P,P)/P$。
于是可把式 (\ref{QTMD}) 改写为
\begin{align}
R(k_3, P)  =  & \,\int_{-\infty}^{\infty} d  k_1
\int_{-1} ^ {1} d  x \,   {\cal F} (x, k_1^2+(k_3-xP)^2 )
\
\label{RTMD0}
\end{align}
或者换用线性宗量 \mbox{$k_3-xP$：}
\begin{align}
R(k_3, P)  =  &  \int_{-1}^1 dx\,   {\cal R} (x, k_3-xP)
\  ,
\label{RTMD}
\end{align}
其中
\begin{align}
{\cal R} (x, k_3  )   \equiv   & \,\int_{-\infty}^{\infty} d  k_1
{\cal F} (x, k_1^2+k_3^2)
\
\label{calD}
\end{align}
是 TMD $ {\cal F} (x, \kappa^2)$ 对二维矢量
$\kappa=\{k_1,k_3\}$ 的 $k_1$ 分量积分后的结果。
根据 (\ref{calD})，${\cal R} (x, k_3  ) $ 通过 $k_3^2$ 依赖于 $k_3$。

对静止强子，
\begin{align}
R(k_3, P=0)  \equiv   &\, r(k_3)=
\int_{-1}^1 dx\,   {\cal R} (x, k_3  )
\  .
\label{DTMDsmall}
\end{align}

这个一维分布可以通过对密度
\begin{align}
\rho (z_3^2) \equiv & {\cal M} (0,z_3^2) =
\int_{-\infty}^{\infty}   dk_3 \,
r(k_3)  \,  e^{i k_3 z_3 } \,
\
\label{dk3}
\end{align}
作参数化直接求得，该密度由
$  \langle p |   \phi(0) \phi ( z_3 )|p \rangle |_{ {\bf p} = {\bf 0} } $ 给出。
因此，$r(k_3)$ 描述静止系强子中 $k_3$
（或 ${\bf k}$ 的任一分量）的原始分布。

公式 (\ref{RTMD}) 有一个直观的解释。
按照它，当强子运动时，部分子的 $k_3$ 动量
有两个来源。

第一部分 $xP$ 来自强子整体的
运动，获得 $xP$ 的概率由 TMD
${\cal F} (x, \kappa^2)$ 对其第一个宗量 $x$ 的依赖性支配。

另一方面，获得其余部分 $k_3-xP$ 的概率
由 TMD 对其第二个宗量 $\kappa^2$ 的依赖性支配，
后者描述了静止系中的原始
动量分布。

参数 $x$ 同时出现在式 (\ref{RTMD}) 中
${\cal R} (x, k_3-xP)$ 的两个宗量里，
也就是说，$R(k,P)$ 是以卷积形式给出的。
在这个意义上，动量分布 $R(k,P)$
以及 quasi-PDFs 都具有混合结构，
同时受到 PDFs 形状与
静止系分布形状的影响。

\subsection{因子化模型}

既然 $k_3$ 的两个来源看起来相互独立，那么用一个因子化模型
\mbox{${\cal R} (x, k_3-xP) = f(x) r (k_3-xP)$}
（$f(x)$ 的 $x$ 积分归一化为 \mbox{1）}
来演示动量分布与 quasi-PDFs 的混合性质是很自然的。
对于原始的 ${\cal M} (\nu,-z^2)$ 函数，这一假设对应于
因子化假定 ${\cal M} (\nu,-z^2) = {\cal M} (\nu,0){\cal M} (0,-z^2)$。

作为演示，我们对静止系密度取 Gaussian 型
$
\rho_G(z_3^2)=  e^{-z_3^2\Lambda^2/4}
$。它对应于
\begin{align}
r_G(k_3) =  &
\frac{1}{\sqrt{\pi } \Lambda}  e^{-k_3^2/\Lambda^2}
\  .
\label{DTMDsmallfac}
\end{align}
对于 $f(x)$，我们取
一个类似核子价密度的简单 PDF：
$f(x)=4(1-x)^3 \theta (0\leq x \leq 1)$。
从图 \ref{R} 可以看到，$R(k,P)$ 曲线
在小 $P$ 时呈 Gaussian 形状，在大 $P$ 时变为类似拉伸 PDF 的形状。

这一结果与一个已知事实完全相符：
运动强子的波函数并不是静止系波函数经过单纯运动学 ``boost''
得到的。
事实上，随着 $P$ 增大，
静止系分布 $r(k)$
的影响越来越弱，最终 $R(k,P)$ 的形状
由一个完全不同的函数 $f(k/P)$ 决定。

换算到变量 $y=k/P$
即得图 \ref{Qgy} 所示的 quasi-PDF $Q(y,P)$。
在大的 $P$ 下，它明显趋于 $f(y)$ 的 PDF 形状。
特别地，取动量
$P \sim 10 \Lambda$ 就能得到相当接近 $P \to \infty$ 极限形状的 quasi-PDF。
不过，由于 $\Lambda \sim \langle k_\perp \rangle$，
按惯用值
$\langle k_\perp \rangle \sim $ 300 MeV 估算，
\mbox{$P\sim 10 \Lambda$} 折合 $P \sim 3$ GeV，
这大得令人不安。于是自然产生一个问题：如何改进收敛性。

\begin{figure}[t]
\centerline{\includegraphics[width=3in]{images/RkP1_10_50keyn.pdf}}
\vspace{-5mm}
\caption{因子化 Gaussian 模型中的动量分布 $R(k,P)$，取 {$P/\Lambda =1,10,50$}。}
\label{R}
\end{figure}

\begin{figure}[t]
\centerline{\includegraphics[width=3.1in]{images/QyP1_10_50key.pdf}}
\vspace{-0.7cm}
\caption{因子化 Gaussian 模型中 quasi-PDF $Q(y,P)$ 的演化，取 {$P/\Lambda =1,10,50$}。
\label{Qgy}}
\end{figure}

\subsection{伪部分子分布（Pseudo-PDFs）}

quasi-PDF
$Q(y,P)$ 结构复杂的
形式上的原因在于：它是把 ${\cal M} (\nu, z_3^2)e^{i \nu y}$
沿 $(\nu, z_3)$ 平面内非水平直线 $z_3=\nu/P$
作 $\nu$ 积分得到的（见式 (\ref{IxM})）。随着 $P$ 增大，斜率
减小，直线趋于水平，
quasi-PDFs 转变为 PDFs。

与此相反，pseudo-PDFs ${\cal P} (x,z_3^2)$
按定义是把 ${\cal M} (\nu, z_3^2)e^{i \nu x}$ 沿水平线
$z_3=$ const 积分得到的。pseudo-PDFs 一个非常吸引人的性质是：
对所有 $z_3$ 值它们都具有
$-1\leq x \leq 1$ 的支撑。对小 $z_3$，
它们转变为 PDFs。

更确切地说，当 $z_3$ 很小时，$1/z_3$
类似于标准 OPE 方法中标度依赖 PDFs $f(x,\mu^2)$ 的重正化参数 $\mu$。

然而有一个微妙之处：PDFs $f(x,\mu^2)$ 的 \mbox{$\mu^2$ 依赖性}
只来自演化对数 $\ln (\mu^2/m^2)$，
而 QCD 中 quasi-PDFs 的 $z_3^2$ 依赖性既来自
演化对数 $\ln (z_3^2 m^2)$，
也来自紫外对数 $\ln (z_3^2 \mu_R^2)$，
其中 $\mu_R$ 是与规范链接重正化相关的发散的截断参数
（见文献 \cite{Craigie:1980qs}）。
在领头对数层面，这些发散不依赖 $\nu$。
结果，``约化'' Ioffe 时间分布
\begin{align}
{\mathfrak M} (\nu, z_3^2) \equiv \frac{ {\cal M} (\nu, z_3^2)}{{\cal M} (0, z_3^2)} \
\label{redm}
\end{align}
在小 $z_3$ 下满足关于 $1/z_3$ 的
领头阶演化方程
\begin{align}
\frac{d}{d \ln z_3^2} \,
{\mathfrak M} (\nu, z_3^2)    = - \frac{\alpha_s}{2\pi} \, C_F
\int_0^1  du \,   B ( u )  \,   {\mathfrak  M} (u \nu, z_3^2)
\label{74}
\end{align}
它与 $f (x,\mu^2)$ 关于 $\mu$ 的演化方程类似。
非单态夸克情形的领头阶演化核 $B(u)$ 由文献 \cite{Braun:1994jq}
给出：
\begin{align}
B (u)    &=  \left [ \frac{1+u^2}{1- u} \right ]_+  \
,
\label{74a}
\end{align}
其中 $[ \ldots ]_+$ 表示标准的 ``plus'' 约定。

对上文所用的模型（并按非单态 PDFs 要求作 $x \to -x$ 对称化），
我们有
${\cal M} (\nu, 0) ={12}\left [\nu^2 - 4\sin ^2( \nu/2) \right ] /{\nu^4}    \  .
$
该函数以及卷积积分
$B \otimes {\cal M}(\nu)$ 的形状示于 \mbox{图 \ref{MBM}。}
可以看到，$B \otimes {\cal M}(\nu)$ 在
$\nu=0$ 处为零，这反映了矢量流的守恒。
因此静止系密度 ${\mathfrak M} (0,z_3^2)$ 不受微扰演化的影响。

\begin{figure}[t]
\centerline{\includegraphics[width=3.3in]{images/MBMkey.pdf}}
\vspace{-0.6cm}
\caption{模型 Ioffe 时间分布 ${\cal M} (\nu,0)$ 以及支配其演化的函数 $B \otimes {\cal M}$。
\label{MBM}}
\end{figure}

\subsection{格点实现}

在格点上求 Ioffe 时间分布的一种可能途径
（由 K.~Orginos 建议）
是对若干个 $P$ 值计算 ${\cal M} (Pz_3, z_3^2)$，
然后用 $\nu$ 和 $z_3^2$ 的函数拟合所得结果。

回想前面关于运动强子 $k_3$ 两个表观独立来源的讨论，
可以期望
${\cal M} (\nu, z_3^2)$ 会因子化，即
${\cal M} (\nu, z_3^2)= {\cal M} (\nu, 0) {\cal M} (0, z_3^2)$。这时
由式 (\ref{redm}) 定义的约化函数 ${\mathfrak M} (\nu, z_3^2) $
等于 $ {\cal M} (\nu, 0)$，
于是求 ${\cal M} (\nu, 0)$ 的目标就达到了。
形式上，剩下要做的只是对其作 Fourier 变换以求得 PDF $f(x)$。

事实上，这种因子化多年前已在
格点 QCD 横动量分布的开创性研究 \cite{Musch:2010ka}
中被观察到。

quasi-PDFs 的一个严重缺点在于：即使在 TMD
[以及 ${\cal M} (\nu, z_3^2)$] 因子化的有利情形下，
它们仍具有 $x$ 卷积结构 (\ref{QTMD})。
相反，使用比值 ${\mathfrak M} (\nu, z_3^2)$ 形式的 pseudo-PDFs，
可以把原始分布的 $z_3^2$ 依赖性约去，
而不影响决定 PDF 形状的 $\nu$ 依赖性。

使用该比值的另一个优点（由 K.~Orginos 指出）是
规范链接 ${ \hat E} (0,z_3; A)$ 的格点重正化所产生的
\mbox{$z_3$ 依赖性}
会相消。
这种重正化的必要性来自线性发散 $|z_3| \delta m$（其中 $\delta m \sim 1/a$，
$a$ 为紫外截断）
与对数发散 $\ln (z_3^2/a^2)$ \cite{Polyakov:1980ca,Dotsenko:1979wb}。
由于其定域性质，预期它们会合并成一个不依赖 $\nu$ 的因子
$Z(z_3^2/a^2)$，
它在比值 ${\mathfrak M} (\nu, z_3^2) $ 的分子和分母中相同。

最近，文献 \cite{Ishikawa:2016znu,Chen:2016fxx}
论证了 ${\cal M} (\nu, z_3^2)$ 线性发散
到所有阶的可乘重正化性。
文献 \cite{Ishikawa:2017faj}
基于对相关 Feynman 图的直接分析，
声称给出了对线性和对数发散的一般证明。

另一种方案 \cite{Ji:2017oey,Green:2017xeu}
是把 $ \hat E(0,z;A) $ 处理为
$ h(0) \bar h (z) $，
其中辅助场 $h(z)$ 类似于重夸克有效理论（HQET）中
无限重夸克场。
由于 HQET 是可乘重正化的 \cite{Bagan:1993zv}，
这意味着 $\bar \psi(0)  \hat E(0,z;A) \psi (z)$ 在微扰论所有阶
也是可乘重正化的。

实际上，${\mathfrak M} (\nu, z_3^2)$ 将带有
残余 \mbox{$z_3^2$ 依赖性。}
它一方面来自软部分可能的因子化破坏
（据文献 \cite{Musch:2010ka} 的结果，
预计相当轻微），
另一方面来自不可避免的微扰演化。对非零 $\nu$，
后者在小 $z_3^2$ 处应表现为
$\ln (1/z_3^2 \Lambda^2)$ 尖峰。

因此，一种建议的策略是从不太小的 $z_3^2$
值（例如大于 0.5 fm$^2$ 的那些）出发，
把 ${\mathfrak M} (\nu, z_3^2)$ 外插到 $z_3^2=0$。
所得函数 ${\cal M} ^{\rm soft}(\nu,0)$ 可视为
给出 ``低归一化点处'' PDF $f_0(x)$
的 Ioffe 时间分布。小 $z_3$ 处剩余的
$\ln (1/z_3^2 \Lambda^2)$ 尖峰将生成它的演化。

要把 ${\mathfrak M} (\nu, z_3^2)$ 转换成 $x$ 的函数，原则上需要知道
所有 $\nu$ 处的 ${\mathfrak M} (\nu, z_3^2)$，这是不可能的。
现有格点计算达到的最大 $\nu$ 值
从 $3\pi $ \cite{Lin:2014zya} 到 $5\pi$ \cite{Alexandrou:2015rja} 及 $6 \pi $ \cite{Orginos:2017kos} 不等。
在这些有限区间上作 Fourier 变换会在 $x$ 中产生非物理振荡。
因此，想法是避开对 $\nu$ 的 Fourier 变换，
直接把从格点得到的约化 Ioffe 时间分布
与从实验已知的部分子分布导出的结果比较。

当然，这一方案的具体技术实现
应留待 ${\mathfrak M} (\nu, z_3^2)$ 的格点数据可用之后再行讨论。
